\chapter{Abdominal Pain in Children}

\section*{Definition and Overview}
Abdominal pain in children refers to discomfort or distress felt in the belly of an infant, child, or adolescent. It is among the most common reasons families present to general practitioners and emergency departments. The overwhelming majority of causes are benign and self-limiting, but a small number are surgical, infective, or otherwise serious, and the range of likely causes shifts markedly with age. Because young children cannot describe their symptoms accurately, clinicians and carers must rely on careful observation of how the child looks, behaves, feeds, sleeps, and passes urine and stool.

\section*{Epidemiology}
Abdominal pain is extremely common throughout childhood. Functional abdominal pain, in which no organic cause is found, affects roughly 1 in 10 school-aged children and is more prevalent than any single structural cause. Constipation and infant colic are very common in the first years of life, and gastroenteritis causes frequent brief episodes at all ages. Appendicitis is uncommon under the age of 5 and peaks in older children and adolescents. Intussusception, a form of bowel obstruction in which one segment of intestine telescopes into another, occurs most often between 6 and 18 months. Urinary tract infections are an important cause of abdominal pain, particularly in infants and young girls.

\section*{Aetiology and Pathophysiology}
The causes of childhood abdominal pain are best considered by age group, since the anatomy, physiology, and exposure of the gut differ across development.

\begin{itemize}
    \item \textbf{Infants:} infant colic, gastroenteritis, constipation, urinary tract infection, cow's milk protein allergy, and gastro-oesophageal reflux
    \item \textbf{Bowel obstruction and torsion:} intussusception, malrotation with volvulus, and incarcerated hernias can obstruct or strangulate the bowel
    \item \textbf{Young children:} gastroenteritis, mesenteric adenitis, constipation, and urinary infection
    \item \textbf{Older children and adolescents:} appendicitis, inflammatory bowel disease, constipation, and functional abdominal pain
    \item \textbf{Adolescent girls:} dysmenorrhoea, ovarian cyst complications, pelvic inflammatory disease, and rarely ectopic pregnancy
    \item \textbf{Boys:} testicular torsion, which can present with abdominal and lower pelvic pain
    \item \textbf{Functional causes:} visceral hypersensitivity and disordered gut-brain interaction, often triggered or aggravated by stress, anxiety, or previous infection
    \item \textbf{Referred pain:} extra-abdominal sources such as lower-lobe pneumonia may be perceived as abdominal pain
\end{itemize}

\section*{Clinical Features}
Assessment begins with the pattern of pain, but in younger children behaviour is often more informative than words.

\begin{itemize}
    \item Site of pain, and whether the child is able to point to a specific spot
    \item In babies, inconsolable crying, drawing the knees up to the chest, or poor feeding
    \item Vomiting, particularly green (bile-stained) or blood-stained vomit
    \item Blood in the stool, which in intussusception classically resembles redcurrant jelly
    \item Fever, diarrhoea, or constipation
    \item Reduced intake of fluids or food, and reduced urine output
    \item Abdominal distension, tenderness, or a rigid abdomen on examination
    \item Pain or burning with urination, suggesting urinary tract infection
    \item In older children, pain that interrupts sleep, school, or play
\end{itemize}

\section*{Differential Diagnosis}
\begin{itemize}
    \item Gastroenteritis versus appendicitis versus mesenteric adenitis in school-aged children
    \item Intussusception versus gastroenteritis in infants and toddlers
    \item Constipation versus functional abdominal pain in young children
    \item Urinary tract infection versus other abdominal causes at any age
    \item Testicular torsion versus other causes of lower abdominal pain in boys
    \item Gynaecological causes versus appendicitis in adolescent girls
    \item Organic disease versus functional pain or anxiety-related symptoms
\end{itemize}

\section*{When to Seek Medical Care}
Parents should seek urgent review if a child has bile-stained vomiting, blood in the stool, a swollen or tender abdomen, a high fever, inability to keep down fluids, or pain severe enough to interfere with normal activity.

\textbf{Call 000 (or local emergency services) for:}
\begin{itemize}
    \item A child who looks seriously unwell: pale, floppy, drowsy, or not responding normally
    \item Difficulty breathing or abnormally rapid breathing
    \item A rigid, tense, or very tender abdomen, or green vomit
    \item Blood in the stool combined with severe crying and drawing up of the legs
    \item Sudden, severe pain with collapse or fainting
\end{itemize}

\section*{Diagnosis and Investigations}
Diagnosis relies on a structured history from the parent or carer, close observation, and targeted tests.

\begin{itemize}
    \item \textbf{History:} onset and site of pain, feeding and drinking, urine output, bowel habit, fever, and any change in behaviour or activity
    \item \textbf{Examination:} vital signs, hydration status, growth, and gentle abdominal palpation, with examination of the groins in boys
    \item \textbf{Urine testing:} dipstick and culture to exclude urinary tract infection
    \item \textbf{Blood tests:} full blood count and inflammatory markers where infection or appendicitis is suspected
    \item \textbf{Ultrasound:} the key investigation for suspected intussusception and helpful in appendicitis and ovarian pathology
    \item \textbf{Plain X-ray or CT:} used where bowel obstruction, perforation, or other serious pathology is possible
    \item \textbf{Family and psychosocial assessment} where functional abdominal pain is suspected
\end{itemize}

\section*{Management}
Management is directed at the underlying cause, with supportive care while the diagnosis is being clarified.

\begin{itemize}
    \item \textbf{Supportive care:} adequate fluids, rest, and paracetamol or ibuprofen at weight-appropriate doses
    \item \textbf{Gastroenteritis:} oral rehydration solution and continuation of normal feeding
    \item \textbf{Constipation:} increased dietary fibre, fluids, and laxatives under medical supervision
    \item \textbf{Appendicitis:} prompt surgical removal of the appendix
    \item \textbf{Intussusception:} pneumatic or hydrostatic reduction under imaging, with surgery reserved for failure or complications
    \item \textbf{Urinary tract infection:} appropriate antibiotics and follow-up where recurrent
    \item \textbf{Functional pain:} reassurance, a regular routine, and psychological support such as cognitive behavioural therapy; medication is rarely required
    \item \textbf{Follow-up:} review to ensure symptoms settle and that growth, sleep, and school attendance return to normal
\end{itemize}

\section*{Complications}
\begin{itemize}
    \item Dehydration from vomiting and reduced fluid intake
    \item Perforation and peritonitis if appendicitis or bowel obstruction is missed
    \item Recurrence of intussusception, particularly after non-surgical reduction
    \item Loss of a testis if torsion is not treated as an emergency
    \item Ongoing functional pain, school absence, and anxiety if the condition is not managed well
    \item Kidney scarring from recurrent or unrecognised urinary tract infections
\end{itemize}

\section*{Prognosis and Outcomes}
Most episodes of abdominal pain in childhood resolve quickly and completely. Appendicitis treated early has an excellent prognosis, and intussusception is usually cured by reduction. Functional abdominal pain frequently improves over months with reassurance and consistent support, although a minority of children continue to experience symptoms into adolescence. The key determinant of outcome is early recognition of serious causes, so prompt medical review of any child who looks unwell is strongly encouraged.

\section*{Prevention and Self-Care}
\begin{itemize}
    \item Encourage thorough hand washing and safe food handling to reduce gastroenteritis
    \item Provide a fibre-rich diet and adequate water to prevent constipation
    \item Maintain regular routines for meals, sleep, and toileting
    \item Keep immunisations up to date, as they protect against several gut infections
    \item Watch for changes in toilet or feeding habits and raise concerns with a doctor
    \item Support children's emotional wellbeing, since stress and anxiety can worsen tummy pain
\end{itemize}

\section*{Sources and Further Reading}
The following trusted organisations provide further reading on childhood abdominal pain for parents, students, and clinicians.

\begin{itemize}
    \item NHS. \textit{Stomach ache in children: causes and when to get help.} https://www.nhs.uk/conditions/stomach-ache/
    \item Mayo Clinic. \textit{Abdominal pain in children: causes and when to see a doctor.} https://www.mayoclinic.org/diseases-conditions/abdominal-pain/
    \item MedlinePlus. \textit{Abdominal pain in children.} https://medlineplus.gov/ency/article/003120.htm
    \item MSD Manual. \textit{Abdominal pain in children: professional reference.} https://www.msdmanuals.com/
    \item Royal Children's Hospital Melbourne. \textit{Clinical practice guidelines on abdominal pain.} https://www.rch.org.au/clinicalguide/
    \item Centers for Disease Control and Prevention. \textit{Childhood health and development resources.} https://www.cdc.gov/
\end{itemize}
